% ================================================================
%  chapters/parts/ch12-c.tex
%  第 12 章  扫描探针显微镜与纳米光谱技术（分片 c）
%  源：S2（下册）第 12 章，书内 p.926—942 / PDF p.412—428
%  范围：12.6（扫描探针与光谱技术的结合：12.6.1 近场扫描光学显微镜、
%        12.6.2 TERS、12.6.3 Nano-FTIR、12.6.4 光学激发的针尖探测方法）
%        与 12.7（结束语）
%  不含：12.1—12.2（STM，分片 a）、12.3—12.5（AFM，分片 b）
%  说明：本分片不写 \chapterstart（分片 a 已写）。
% ================================================================

\subsection{扫描探针与光谱技术的结合}
\label{sec:nano-intro}

\begin{probout}
扫描探针显微镜把空间分辨率推到了原子尺度，但常规 STM 与 AFM 给不出化学信息：表面元素组成、
吸附物种的种类与吸附构型都不在它们的输出里。本节要回答的是——\textbf{怎样在保留纳米乃至
亚纳米空间分辨率的同时，拿到分子的振动（化学）指纹}，以及为此付出的代价是什么。原始输出
有三类：逐点采集的\textbf{近场拉曼谱或近场红外谱}（峰位以 cm$^{-1}$ 计）、由某个特征峰的
强度或峰面积逐点拼出的\textbf{化学像}、以及与之同时记录的\textbf{形貌像}；此外还有一组必须
随图报告的条件（针尖材质与曲率半径、波长与功率、针尖—样品间距、解调阶数、每像素采谱时间、
环境）。本节的核心结论只有一句：\textbf{空间分辨率不再由波长（衍射极限）决定，而由针尖顶端
的曲率半径与近场增强体积决定；代价是信号极弱、探针易污染钝化、扫描极慢，并且远场背底必须被
有效抑制}。后面多种技术的差别，本质上只是这四项代价的不同付法。
\end{probout}

\paragraph{化学信息的缺口}
STM 与 AFM 之外，高空间分辨的技术在化学识别上各有短板。STM 路线上的非弹性电子隧道谱
（inelastic electron tunneling spectroscopy，IETS）已能提供原子分辨率的表面振动谱，达到了
目前分子振动谱学的空间分辨率极限；但隧穿电子的振动吸收截面极低，受信号灵敏度与隧穿结构
稳定性的限制，IETS 目前只能在极低温（通常为液氦温度）与超高真空下工作，且能探测的振动模式
相当有限\src{S2 §12.6，书内 p.926 / PDF p.412}。透射电子显微镜虽然能在原子尺度给出固体
的原子与电子结构，却常常只能对重元素成像，因此在催化表征中往往只能对纳米粒子本身、而不是
吸附分子做高分辨成像。而"表面分子与反应中间体到底是什么"恰恰是红外、拉曼等分子振动光谱
所擅长的——问题在于，这些经典光谱技术受光学衍射极限限制，其宏观平均信息在解释局域活性位
特征时经常失灵\src{S2 §12.6，书内 p.926—927 / PDF p.412—413}。

\paragraph{衍射极限：全部问题的出发点}
\begin{principle}
传统远场光学成像的空间分辨率受光波衍射限制。点光源经透镜成像时呈现艾里斑，19 世纪末
Abbe 与 Rayleigh 由此推出瑞利判据：两个点源可被清楚分辨的最小距离为
\begin{equation}\label{eq:rayleigh}
\Delta x \;\geq\; \frac{0.61\,\lambda}{n\sin\theta}
\end{equation}
式中 $\lambda$ 为照明光波长，$n\sin\theta$ 为数值孔径（numerical aperture，N.A.）。由于通常
$n<2$、$\sin\theta<1$，可分辨的最小距离一般不小于 $\lambda/2$，即成像波长的一半
\src{S2 §12.6，书内 p.927 / PDF p.413}。
\emph{核对说明}：已对照原页面图像（书内 p.927 / PDF p.413）。原书印「两个点源之间可被清楚
分辨的最小距离 $\Delta x \geq 0.61\lambda/n\sin\theta$，其中，$\lambda$ 为照明光波长，
$n\sin\theta$ 为数值孔径（numerical aperture，N.A.）」；「由于通常 $n<2$、$\sin\theta<1$，
所以可分辨的最小距离 $\Delta x$ 一般不小于 $\lambda/2$，即成像光波长的一半」；「其在可见光下
的最大空间分辨率是 $200\sim400$ nm，如果使用中红外光谱，则空间分辨率要减至几微米的量级」——
数值与变量定义均与本分片一致。原书正文是行内线性写法（OCR 亦作线性串），本分片按通行写法排
成分式，系数 $0.61$ 照录。
\end{principle}

该判据给出现成的数量级：传统光学系统在可见光下的最大空间分辨率是 200—400 nm；若使用中红外
光谱，空间分辨率要降到几微米量级\src{S2 §12.6，书内 p.927 / PDF p.413}。这正是"红外光谱
擅长化学识别、却不擅长纳米成像"的根源——红外波长本身（常用 2.5—25 µm，对应 400—
4000 cm$^{-1}$）已经从根本上限制了它的空间分辨能力。

绕开衍射极限有几条路，各自代价不同。其一是换更短的波长，例如高能电子与离子，波长可小于
一个原子，从而实现原子分辨率成像；但很高的能量通常带来严重样品损伤。其二是荧光标记：用染料
分子做化学标记，在可见光区激发、由激发态弛豫发出荧光，配合合适的滤光片与探测器可实现单分子
荧光成像；共聚焦荧光显微镜结合先进光源与探测器后，已发展出受激发射损耗显微术（STED）、
光活化定位显微术（PALM）、随机光学重构显微术（STORM）等方法，可获得远低于衍射极限的空间
分辨率。这些超分辨远场技术在生物研究中应用广泛，但\textbf{它们能提供的光谱信息仍相当有限}
\src{S2 §12.6，书内 p.927 / PDF p.413}。其三是本节的主题：把探针放进光路，用近场代替远场。
事实上在 STM 与 AFM 发明后不久，就有人尝试把扫描探针的高空间分辨与光谱的化学分辨结合起来；
目前近场光学显微镜、针尖增强光谱技术以及光学激发的针尖探测方法，都被认为具有用于催化原位
研究的潜力\src{S2 §12.6，书内 p.927 / PDF p.413}。

\subsubsection{近场扫描光学显微镜}
\label{ssec:nsom}

\paragraph{从 Synge 的小孔到第一台近场显微镜}
1928 年 Synge 提出用小孔径实现超衍射极限光学分辨率的设想：光通过小孔后高度发散，在远场中
当传播距离远大于波长与孔径时，衍射导致强度起伏、形成艾里斑；但在近场区域，只要物体足够
接近小孔，成像分辨率就能接近孔径尺寸、远小于波长\src{S2 §12.6.1，书内 p.927—928 /
PDF p.413—414}。由于当时难以对纳米尺度的间距做精确控制，这一可能性直到 1972 年才被实验
证实。STM 实现了探针对表面的快速精确扫描与针尖—样品间隙的精确控制，直接推动了近场扫描
光学显微镜（near-field scanning optical microscopy，NSOM，也称 SNOM）的开发：1984 年，
在瑞士 IBM 研究实验室工作的 Pohl 等制造了第一台用可见光测量的近场光学显微镜，通过探针在
样品表面保持数十纳米的距离采集反馈信息，并在几年后实现高分辨成像
\src{S2 §12.6.1，书内 p.928 / PDF p.414}。

\paragraph{孔径式近场显微镜：原理与它自己的天花板}
最初的方法是用金属（如 Al 或 Ag）涂覆的锥形光纤作探头（光纤底部孔径约 100 nm），借助扫描
探针的反馈机制把光纤靠近样品表面，光通过底部小孔照射样品，从而避开衍射形成的艾里斑，在
测量形貌的同时以远低于波长的分辨率成像——这就是\textbf{孔径式扫描近场显微镜}。采集方式随
样品而变：样品对探测光透明时，可在样品下方检测透射光子，或在周围检测激发的荧光；不透明
样品则检测反射光；也可以从外部照射样品、再用光纤收集光子送到探测器。其中特别值得注意的是
\textbf{光子扫描隧道显微镜}（PSTM，亦称扫描隧道光学显微镜 STOM）：光从样品背面全内反射产生
隐失波（evanescent wave），光纤探针端部检测样品表面的隐失场并把它转换为传播模式送到远处
的探测器；隐失波受样品调制，因此近场光学像携带局域样品信息，但图像解析通常较复杂，源书
称目前尚未见用于催化研究的报道\src{S2 §12.6.1，书内 p.928 / PDF p.414}。

孔径式的主要弱点在于小孔径光纤的低透过率。为补偿这一点，光源通常都是高亮度、低发散的激光；
即便如此，对 80—100 nm 的孔径，光纤传导的可见光透过率也只有 $10^{-7}$—$10^{-4}$ 量级，
即几十纳瓦的功率，空间分辨率因此严重受限。更麻烦的是透过率随孔径减小或波长变长而进一步
下降：中红外波段下 50 nm 孔径只剩 $10^{-30}$ 量级，红外近场显微镜用孔径式方案几乎不可行
\src{S2 §12.6.1，书内 p.928 / PDF p.414}。

\paragraph{散射式近场显微镜：把针尖当作纳米散射体}
为克服孔径式的限制，人们发展了散射式近场光学显微镜（scattering-type SNOM，s-SNOM）：用
激光聚焦照射扫描探针的金属针尖，在针尖附近激发一个纳米尺度的增强近场区域；针尖接近样品
表面时，由于不同物质介电性质不同，近场光学信息被改变；通过精确调整实现\textbf{焦斑—针尖—
样品三者对准}，散射光经显微物镜由远场探测器收集，再经精密的调制过滤技术解析，就能得到表面
的近场光学谱并成像\src{S2 §12.6.1，书内 p.928—929 / PDF p.414—415}。以金属针尖顶点作为
纳米散射体，s-SNOM 突破了孔径式的限制：\textbf{其空间分辨率与光学分辨率通常被认为仅由探针
针尖的曲率半径决定}，使几乎所有的光谱技术都能扩展到近场使用，成像分辨率不再受波长限制，因而
尤其适合红外与太赫兹波的微观成像，并可能具有 10 nm 或更好的空间分辨率
\src{S2 §12.6.1，书内 p.929 / PDF p.415}。

\paragraph{分辨率的来源与四项代价（本节核心）}
\begin{principle}
\textbf{分辨率来源。}近场方案把"分辨"这件事从波长转移到了探针几何上。近场（隐失）信号随
离开表面的距离迅速衰减，可用单一衰减长度 $\delta$ 描述：
\begin{equation}\label{eq:nearfield-decay}
E_{\mathrm{nf}}(z)=E_{\mathrm{nf}}(0)\,e^{-z/\delta},\qquad \delta \sim a
\end{equation}
式中 $z$ 为针尖—样品间距，$a$ 为针尖顶端曲率半径。只要间距 $z$ 与针尖曲率半径同量级，
参与相互作用的样品体积就被限制在针尖正下方，其横向尺寸由 $a$（而非 $\lambda$）决定，这就是
"分辨率由针尖曲率半径决定"的物理含义。
\emph{核对说明}：已对照原页面图像（书内 p.928—929 / PDF p.414—415）。原书对近场衰减只作
文字说明：光通过小孔在远场高度发散，而在近场区域物体足够接近小孔即可获得接近孔径尺寸、远小于
光波长的分辨率（书内 p.928）；「其空间分辨率和光学分辨率通常认为仅由扫描探针针尖的曲率半径
决定」（书内 p.929）——原书\textbf{未给出}解析式，也\textbf{未定义} $\delta$。式
~\eqref{eq:nearfield-decay} 按隐失场的标准形式补写，$\delta\sim a$ 为量级估计（隐失场衰减长度
与针尖曲率半径同量级），非原书原文。另需注意：「分辨率由针尖曲率半径决定」一语在原书中出现
于 §12.6.1 末段对散射式近场显微镜（s-SNOM）的论述；孔径式近场显微镜的分辨率则由孔径尺寸决定
（书内 p.928）。

\textbf{代价一：信号极弱。}只有针尖正下方极小的体积贡献近场信号，且信号随间距指数衰减
（式~\eqref{eq:nearfield-decay}）。这在数值上是残酷的：孔径式方案中中红外 50 nm 孔径的透过率
仅 $10^{-30}$ 量级；散射式方案中，从悬臂与针尖杆散射出来的\textbf{远场背底远强于针尖末端
散射的近场信号}。因此近场光谱的可行性从来不取决于"有没有信号"，而取决于\textbf{能不能把
极弱的近场信号从强得多的远场背底里分离出来}。

\textbf{代价二：探针易污染、易钝化。}增强场与分辨率都由针尖最末端几纳米的几何与光学性质
决定，任何吸附物、氧化层或磨损都会同时改变二者；针尖还可能与样品发生化学作用而把自身变成
催化剂（见 \ref{ssec:ters} 段的 Ag 针尖与氧化铝保护层）。

\textbf{代价三：扫描极慢。}近场信号的绝对强度低，每个点位都必须积分足够长的时间才能得到可用
谱图。源书给出的两个标尺是：同步辐射红外光源下单点谱约几秒、化学成像需几分钟；而 Gross 等
的实验中单点红外谱约需 20 min，以致只能做线扫描而无法成像。

\textbf{代价四：远场背底必须被有效抑制。}这是近场光谱与常规光谱最大的操作差别。抑制手段是
调制—解调：让针尖以频率 $\Omega$ 在 $z$ 方向振动，近场信号因式~\eqref{eq:nearfield-decay}
而被周期性调制，在探测器中产生 $n\Omega$（$n=1,2,3,4,\dots$）的谐波，未受调制的远场背底
（直流部分）则留在低频；锁相放大只提取 $n\Omega$ 分量，就压掉了背底。其数学形式就是傅里叶
分量：
\begin{equation}\label{eq:demod}
S_n=\frac{\Omega}{\pi}\int_{0}^{2\pi/\Omega} I\big(z(t)\big)\cos(n\Omega t)\,\dd t,
\qquad z(t)=\bar z+A\cos\Omega t
\end{equation}
\textbf{这四项代价是同一件事的四个侧面}：因为信号只来自极小体积，所以它弱、所以必须慢、
所以探针状态决定一切、所以必须靠调制把背底和信号分开。判断一台近场实验是否可信，就是逐条
问这四项是否被处理过。
\end{principle}

\begin{figure}[htbp]
\centering
\begin{tikzpicture}[x=1mm,y=1mm,>=Stealth]
% ---- 坐标区间静态核算（x=1mm，y=1mm；版心 165 mm，图宽上限 \onecolwidth = 132 mm）----
% 图元：样品矩形 (6,10)-(98,18)；吸附分子标记 x=40,46,52,58,64，y 18→20.5；
%       针尖多边形 (52,22.5)-(45,36)-(49,36)-(49,52)-(55,52)-(55,36)-(59,36)-cycle；
%       近场热点椭圆 中心(52,20)、rx=9、ry=4 → x[43,61] y[16,24]；
%       z 向振动双箭头 x=42，y 20→30；
%       激发光束两线 (10,70)→(52,20)、(22,70)→(52,20)；
%       近场收集线 (61,20)→(86,58)；远场背底线 (54,50)→(86,66)；
%       探测器盒 (86,56)-(126,68)。
% 节点区间（anchor 已注明；宽度按 text width 计，纵向按注1行±3.5mm、注2行±5.5mm、注3行±6.5mm）：
%  N1 样品  中心(52,14)      w44     → x[30,74]      y[12,16]      在样品矩形内
%  N2 近场区 (63.5,20) west  w26     → x[63.5,89.5]  y[18,22]
%  N3 针尖  (34,25) east     w26     → x[8,34]       y[18.5,31.5]（3 行）
%  N4 z     (40.5,26) east   自然宽   → x[35.5,40.5]  y[24,28]
%  N5 远场背底 (60,48) west  w16     → x[60,76]      y[46,50]
%  N6 近场信号 (73,34) west  w14     → x[73,87]      y[32,36]
%  N7 探测器 中心(106,62)    w34     → x[89,123]     y[58,66]      在探测器盒内
%  N8 激光  (4,72) north west w40    → x[4,44]       y[72,76]
% 两两求交核算（只列有风险的组合，均已排除）：
%  N3∩N4：x [8,34]∩[35.5,40.5]=∅（间隙 1.5 mm）；
%  N3∩激发光束：光束在 y∈[18.5,31.5] 处 x∈[43.2,52.4]，与 [8,34] 不交；
%  N4∩激发光束：光束在 y∈[24,28] 处 x∈[45.3,49.6]，与 [35.5,40.5] 不交；
%  激发光束∩z 双箭头：光束 A 在 x=42 处 y=31.9>30（箭头顶端），不相交；
%  N2∩近场收集线：收集线在 y∈[18,22] 处 x∈[59.7,62.3]，与 [63.5,89.5] 不交；
%  N6∩近场收集线：收集线在 y∈[32,36] 处 x∈[68.9,71.5]，与 [73,87] 不交；
%  N5∩近场收集线：收集线在 y∈[46,50] 处 x∈[78.1,80.7]，与 [60,76] 不交；
%  N5∩远场背底线：背底线在 x∈[60,76] 处 y∈[53,61]，与 [46,50] 不交；
%  N6∩远场背底线：y 32..36 与 59.5..66 不交；N6∩N5：y [32,36]∩[46,50]=∅；
%  N5／N6∩探测器盒：x ≤87 < 86? 否——N6 x 至 87>86，但 y 36<56，故不相交；
%  N8∩激发光束：光束起点 y=70<72；N8∩探测器盒：y 72>68；
%  样品矩形∩N3：N3 y 下界 18.5>18；分子标记 x 40..64 与 N3 x ≤34 不交。
% 最大 x 坐标 = 126（探测器盒右边），最右文字到 x=123；最小 x = 4（N8 左界）。
% 故图宽 ≈ 122 mm < 132 mm（\onecolwidth），图高 ≈ 66 mm。
% ---- 样品与表面吸附分子层 ----
\fill[black!10] (6,10) rectangle (98,18);
\draw[axisline] (6,10) rectangle (98,18);
\node[align=center, text width=44mm] at (52,14) {样品与表面吸附分子层};
\foreach \p in {40,46,52,58,64}{\draw[thinline] (\p,18) -- (\p,20.5);}
% ---- 金属针尖（同时是光学天线）----
\fill[black!25] (52,22.5) -- (45,36) -- (49,36) -- (49,52) -- (55,52) -- (55,36) -- (59,36) -- cycle;
\draw[mainline] (52,22.5) -- (45,36) -- (49,36) -- (49,52) -- (55,52) -- (55,36) -- (59,36) -- cycle;
\node[anchor=east, align=right, text width=26mm] at (34,25) {金属针尖（Ag/Au、PtSi）};
% ---- 近场增强区：横向尺寸由针尖曲率半径决定 ----
\fill[black!8] (52,20) ellipse [x radius=9, y radius=4];
\draw[dashed, accentsoft] (52,20) ellipse [x radius=9, y radius=4];
\node[anchor=west, align=left, text width=26mm] at (63.5,20) {近场增强区};
% ---- 针尖 z 向振动（频率 Ω，锁相解调其 nΩ 分量）----
\draw[<->, accent] (42,20) -- (42,30);
\node[anchor=east] at (40.5,26) {$z$};
% ---- 聚焦激光：焦斑远大于针尖 ----
\draw[thinline] (10,70) -- (52,20);
\draw[thinline] (22,70) -- (52,20);
\node[anchor=north west, align=left, text width=40mm] at (4,72) {聚焦激光（焦斑远大于针尖）};
% ---- 近场信号与远场背底走同一路光 ----
\draw[->, mainline] (61,20) -- (86,58);
\draw[->, mainline] (54,50) -- (86,66);
\node[anchor=west, align=left, text width=14mm] at (73,34) {近场信号};
\node[anchor=west, align=left, text width=16mm] at (60,48) {远场背底};
% ---- 探测器与锁相解调 ----
\fill[black!6] (86,56) rectangle (126,68);
\draw[axisline] (86,56) rectangle (126,68);
\node[align=center, text width=34mm] at (106,62) {探测器 + 锁相解调（$n\Omega$）};
\end{tikzpicture}
\caption{针尖—样品近场耦合与远场背底抑制的几何。读者应看三件事：① 参与相互作用的样品体积
由针尖顶端的\textbf{近场增强区}决定，其横向尺寸与针尖曲率半径同量级，与聚焦焦斑（微米量级）
无关——这是"分辨率不再受衍射极限限制"的几何来源；② 近场信号与来自针尖杆、悬臂的\textbf{远场
背底}走的是同一路散射光，进入同一个探测器，因此必须靠针尖 $z$ 向振动（频率 $\Omega$）调制、
再锁相解调 $n\Omega$ 分量才能把二者分开（式~\eqref{eq:demod}）；③ 近场信号随间距 $z$ 指数
衰减（式~\eqref{eq:nearfield-decay}），所以针尖必须始终停在距表面几纳米处，这也正是探针
污染、磨损与会话式扫描缓慢的根本原因。图为示意，不按真实比例（真实针尖—样品间距约几纳米，
而焦斑为微米量级）。\schemamark\figsrc{据 S2（下册）§12.6.1，书内 p.927—929 / PDF p.413—
415 正文改绘。}}
\label{fig:snom-geometry}
\end{figure}

\begin{figure}[htbp]
\centering
\begin{tikzpicture}
\begin{axis}[
  width=\onecolwidth,
  height=6.2cm,
  xlabel={针尖—样品间距 $z$ / nm},
  ylabel={归一化信号强度},
  xmin=0, xmax=20, ymin=1e-5, ymax=3,
  ymode=log, log basis y=10,
  ytick={1e-4,1e-3,1e-2,1e-1,1},
  yticklabels={$10^{-4}$,$10^{-3}$,$10^{-2}$,$10^{-1}$,$1$},
  xtick={0,5,10,15,20},
  domain=0:20, samples=200,
  legend pos=north east,
  legend cell align=left,
]
\addplot[accent, solid] {exp(-x/2)};
\addlegendentry{近场信号 $\propto e^{-z/\delta}$（$\delta=2$ nm）}
\addplot[black, dashed] {0.01};
\addlegendentry{远场背底（取为接触时近场的 $1\%$）}
\addplot[black!55, dotted, line width=1.4pt] {1e-4};
\addlegendentry{高次解调（$n\Omega$）后的残余背底}
\end{axis}
\end{tikzpicture}
\caption{为什么近场实验"必须靠近、必须解调"——一个只用于说明数量级关系的模型曲线。实线是
式~\eqref{eq:nearfield-decay} 取 $\delta=2$ nm 的近场信号；虚线是远场背底，示意性地取为接触
时近场信号的 $1\%$；点线是经高次解调后残余的背底。读者应看交点：\textbf{若不做任何抑制},
针尖退到约 9 nm 以外，近场信号就落到背底之下；把背底再压低两个数量级，可用的工作距离才
扩展到约 18 nm。结论不是这两个具体数字，而是"可用工作距离由近场信号与残余背底之比决定"。
\modelmark 曲线为模型计算，非实测数据；衰减长度与背底比例均为示意值，源书未给出。
\emph{核对说明}：已对照原页面图像（书内 p.929、p.936 / PDF p.415、p.422）。原书只作定性
说明：「由于从 AFM 探针悬臂和针尖散射出来的远场背景信号远强于从 AFM 针尖末端散射出来的近场
信号，因此必须对近场红外信号的光强进行调制」；「通过锁相放大处理就可提取这些近场信号，同时
剔除未受调制的背景信号（直流部分）」（书内 p.936）。原书\textbf{未给出}曲线与数值：本图的衰减
长度 $\delta=2$ nm、背底取接触时近场信号的 $1\%$、高次解调后再降两个数量级，均为示意取值，
只用于说明量级与"可用工作距离由近场信号与残余背底的比值决定"这一结论。
\figsrc{据 S2（下册）§12.6.1、§12.6.3，书内 p.929、p.936—937 / PDF p.415、p.422—423 的
定性描述绘制。}}
\label{fig:nearfield-decay}
\end{figure}

\begin{measure}
近场实验的条件清单——缺一项，图与结论就不能跨实验比较。\textbf{（1）针尖}：材质（Ag、Au 或
Ag/Au 包覆；AFM 线路上常用 PtSi 作光学天线）、曲率半径、是否加保护层；\textbf{（2）间距与
调制}：针尖—样品间距 $z$、振动振幅 $A$ 与频率 $\Omega$、解调阶数 $n$（近场信号在 $n\Omega$
上的分布与背底完全不同，$n$ 决定了背底被压掉多少）；\textbf{（3）光}：波长（决定能否共振）、
功率密度（过高会引发扩散、脱附甚至分子损坏）、焦斑—针尖—样品是否三者对准；\textbf{（4）时间}：
每像素积分时间与总采谱时间（决定时间分辨率，也决定热漂移累积）；\textbf{（5）环境}：大气、
液相还是超高真空与低温；\textbf{（6）参考}：是否记录了针尖离开表面（如回撤 5 nm）的对照谱。
源书给出的一个具体对照是：图 12-31（c）中黑线即针尖远离表面 5 nm 时采集的 TERS 谱线——没有
这类对照，就无法判断读到的峰究竟来自近场还是远场
\src{S2 §12.6.2，书内 p.933 / PDF p.419}。
\end{measure}

\subsubsection{针尖增强拉曼光谱技术（TERS）}
\label{ssec:ters}

拉曼光谱已被广泛用于原位检测多相催化剂及其表面化学反应，其灵敏度低的问题可由共振拉曼或
表面增强拉曼光谱（SERS）等方式缓解，但成像上仍受光学衍射极限限制，常规远场手段无法突破；
因此对化学反应做纳米尺度监测、尤其是对单粒子催化活性的实时监测，至今仍是显著挑战。2000 年，
美国、日本、德国与瑞士的几个研究组几乎同时把拉曼光谱与特殊设计的扫描探针结合起来，提出了
针尖增强拉曼光谱技术（tip-enhanced Raman spectroscopy，TERS）：用扫描探针的针尖探测底部
样品，同时限域并增强入射光场强，从而同时获得高空间分辨率与高检测灵敏度的拉曼成像
\src{S2 §12.6.2，书内 p.929 / PDF p.415}。TERS 结合了 SERS 的高化学敏感性与 SPM 的高空间
分辨率，是目前发展最迅速、并已实现单分子识别的纳米光谱技术；它已被用于光电器件、碳纳米管、
半导体、生物学与二维材料等领域，但用于催化反应的研究尚处起步阶段，主要是用模型反应展示
其高分辨光谱成像在催化研究中的可行性\src{S2 §12.6.2，书内 p.929 / PDF p.415}。

\begin{principle}
\textbf{增强机制：局域表面等离激元把光场压缩到针尖末端。}把金属（或涂覆薄层 Ag／Au）的 SPM
针尖靠近样品表面并置于激光焦斑中心，适当波长的照射会在针尖末端激发强烈的局域表面等离激元
共振（localized surface plasmon，LSP），增强几纳米距离内的电磁场强。源书给出的电磁增强效率
量级为 $10^{3}$—$10^{9}$，取决于\textbf{针尖的曲率半径、针尖与样品的间距以及所用材料}；正是
这一局域增强使针尖下方吸附分子的拉曼信号急剧提升，从而实现突破衍射限制的空间分辨率
\src{S2 §12.6.2，书内 p.929 / PDF p.415}。若把增强效率写成增强因子，拉曼散射的强度正比于
局域场强的四次方（激发与斯托克斯两个频率各贡献平方项）：
\begin{equation}\label{eq:ters-ef}
\mathrm{EF}\;\approx\;
\left|\frac{E_{\mathrm{loc}}(\omega_{\mathrm{L}})}{E_{0}}\right|^{2}
\left|\frac{E_{\mathrm{loc}}(\omega_{\mathrm{S}})}{E_{0}}\right|^{2}
\;\approx\;\left|\frac{E_{\mathrm{loc}}}{E_{0}}\right|^{4}
\end{equation}
式中 $\omega_{\mathrm{L}}$、$\omega_{\mathrm{S}}$ 分别为激发光与斯托克斯光的频率，$E_0$ 为
无针尖时的入射场。末端的近似在斯托克斯位移远小于激发频率时成立。
\emph{核对说明}：已对照原页面图像（书内 p.929 / PDF p.415）。原书印「理论预测该电磁场的增强
效率在 $10^{3}\sim10^{9}$（取决于针尖的曲率半径、样品和针尖间的距离及所用的材料等）」——即原书
\textbf{只给增强效率的量级与依赖因素，未给增强因子的表达式}，也\textbf{未定义} $E_0$。式
~\eqref{eq:ters-ef} 按电磁增强机制的标准形式补写：$\omega_{\mathrm{L}}$、$\omega_{\mathrm{S}}$
分别为激发与 Stokes 频率，$E_{0}$ 为无针尖时的入射场，当 Stokes 位移远小于激发频率时四次方
近似成立。四阶形式是本分片的写法，非原书原文。
\end{principle}

\paragraph{TERS 与 SERS 的同与不同}
\textbf{相同之处}在于物理机制：两者都靠金属纳米结构的 LSP 局域电磁场增强拉曼散射，增强量级、
对 Ag／Au 的依赖、对激发波长的共振要求也属同一族。\textbf{不同之处有三条，且都指向"可寻址性"}\\
① \emph{热点的可控性不同}：SERS 的增强来自电极粗糙化、胶体聚集或粗糙膜，热点位置随机且
不可预知，测到的是大量分子的系综平均；TERS 只用一根 SPM 针尖作为唯一热点，热点位置由压电
扫描决定，并需精确实现焦斑—针尖—样品三者对准\src{S2 §12.6.1，书内 p.928—929 /
PDF p.414—415}。\\
② \emph{空间编码不同}：TERS 的每个像素有确定坐标，且与同时记录的形貌像严格对应，因而可以把
"一张平均谱"变成"逐点位谱"；SERS 本身没有空间编码能力，其成像分辨率仍受衍射极限与热点分布
限制。\\
③ \emph{采样体积与代价不同}：TERS 的增强体积只有针尖末端几纳米，绝对信号低、对间距极其
敏感、针尖污染或磨损会同时降低增强与分辨率；SERS 的增强体积大得多、信号强，但分子必须贴在
金属表面上，且无法选择"在哪个位置测"。\\
\textbf{为什么 TERS 能做单点位化学成像：}正因为增强场被局域在针尖正下方、并按
式~\eqref{eq:nearfield-decay} 随距离迅速衰减，"针尖停在哪里采谱，谱就只报告哪里的分子"。
把针尖逐点移动、在每个像素上只取某一条特征峰（例如 DMAB 的 1142 cm$^{-1}$）的强度，就得到
化学像；同一次扫描的 STM／AFM 通道提供形貌配准。这也说明 TERS 的分辨率上限来自增强体积
而非波长——在最好的条件下（超高真空、低温、有序的针尖—分子纳米腔），增强体积可以收缩到
亚纳米，源书据此把 TERS 的空间分辨率推到 0.5 nm\src{S2 §12.6.2，书内 p.933—934 /
PDF p.419—420}。

\paragraph{模型反应一：Ag 针尖既是天线又是催化剂}
光催化生成 $p,p'$-二巯基偶氮苯（DMAB）的过程是 TERS 研究中常用的模型反应。Weckhuysen 等
率先用 TERS 原位实时研究了 Au 纳米片上光催化还原对硝基苯硫酚（pNTP）生成 DMAB 的过程
（源书图 12-28）：Ag 包覆的 TERS 针尖受可见光照射时，通过 LSP 共振增强局域电磁场，既增强
了散射的拉曼信号，又有利于光化学中的电荷转移——Ag 针尖接触样品时还作为催化剂在金属表面
产生热电子与空穴，进一步激活吸附分子与催化反应。该工作采用两个激光源的组合：532 nm 绿光
用于激活光催化反应，633 nm 红光仅用于监测、不触发反应（源书图 12-29）。这一设计的意义在于
\textbf{把"自然发生的过程"与"光催化反应"分开}：样品表面的扩散与分子取向变化同样会改变光谱
特征，只有用一个不引发反应的波长采谱，才能把它们与真正的反应产物区分开。作者由此展示了
Ag 针尖末端作为单一催化"粒子"的实时工况研究，并利用金纳米片上的分子自组装展示了 TERS
研究单层分子的高灵敏度\src{S2 §12.6.2，书内 p.929—930 / PDF p.415—416}。源书图 12-29 的
时间序列给出的关键判据是：pNTP 的特征峰位于 1335 cm$^{-1}$、DMAB 的特征峰位于
1440 cm$^{-1}$，100—130 s 的绿光辐照用阴影带示出，两条对照谱线分别取自 90 s 与 265 s
\src{S2 §12.6.2，书内 p.931 / PDF p.417}。

\paragraph{模型反应二：把针尖的催化活性"关掉"}
上述针尖在增强拉曼灵敏度的同时也具催化活性，为使 TERS 更普适，Roy 与 Wain 等在 2015 年用
超薄氧化铝膜保护 Ag 针尖：既阻止针尖与反应物分子发生化学作用，又保持其等离激元增强功能。
用这种保护针尖，他们展示了对纳米银基底上对巯基苯胺（pMA）到 DMAB 的光催化氧化过程的
纳米级成像（源书图 12-30），成像的最佳空间分辨率为 20 nm
\src{S2 §12.6.2，书内 p.930—932 / PDF p.416—418}。这个例子同时给出了本节要强调的边界：
\textbf{20 nm 已远小于激发波长，但仍不足以分辨单个分子}；而且局域等离激元场的空间范围似乎
是空间分辨率的主要限制之一，传统 TERS 通常还需要强入射激光，可能引发扩散、脱附甚至损坏
分子等副反应，影响成像的可持续性与稳定性\src{S2 §12.6.2，书内 p.931—932 /
PDF p.417—418}。

\paragraph{读图实例：TERS 化学成像究竟分辨到了什么（源书图 12-30）}
\begin{obs}
图 12-30 由四个面板组成，图注与正文给出的可观察事实如下。（1）（a）是玻璃衬底上 Ag 纳米粒子
的 AFM 高度图。（2）（b）是对（a）中虚线方框区域、以 DMAB 的 1142 cm$^{-1}$ 吸收峰作出的
空间 TERS 成像。（3）（c）是从（b）所示位点取出的近场 TERS 光谱：位点 1、2、3 的近场光谱
显示出较强的 DMAB 特征 Raman 峰，位点 4、5、6 则以反应物 pMA 的特征峰（1595 cm$^{-1}$）为主。
（4）（d）是沿（b）中一条直线作的强度剖面，由图注给出的空间分辨率为 20 nm
\src{S2 §12.6.2，图 12-30 图注，书内 p.932 / PDF p.418}。
\end{obs}
\begin{infer}
由以上观察可以作出三级推断，并各自带有不同的确定性。\textbf{第一级（几乎确定）}：同一视场内
不同位点的近场光谱归属于不同物种（产物 DMAB 与反应物 pMA），说明光催化转化在纳米尺度上是
\textbf{空间不均匀}的——不是整个基底一起反应，而是部分粒子或部分位点上发生了反应。
\textbf{第二级（确定，但需要条件）}：20 nm 的成像分辨率\textbf{远小于}所用激光波长（可见光
几百纳米），因此这一分辨率不可能来自光学成像系统本身，只能来自针尖末端的近场增强体积——
这正是"分辨率由针尖曲率半径而非波长决定"的直接实验证据（图~\ref{fig:snom-geometry}）。
\textbf{第三级（候选解释，需回查）}：位点 4、5、6 以反应物峰为主，可能是该处未发生反应，也
可能是产物扩散离开、或该处等离激元增强弱使产物峰淹没在噪声里；仅凭这一幅化学像无法区分。
同理，"某类粒子活性更高"的结论还需要把化学像与（a）的粒径、形貌、聚集状态逐粒对照。
至于把 20 nm 的化学像读成"单分子成像"则是过度解读：源书明确说明该分辨率仍不足以实现化学
分辨单个反应物或产物分子。
\emph{核对说明}：已对照原页面图像（书内 p.931—932 / PDF p.417—418）。图 12-30 图注印「（c）为
（b）中所示位点的近场 TERS 光谱图。其中 1、2、3 位点的近场光谱显示出较强的 DMAB 的特征 Raman
峰，而 4、5、6 位点则以反应物 pMA 的特征峰（\qty{1595}{cm^{-1}}）为主。（d）为（b）中所示直线
的强度剖面图，显示出 \qty{20}{nm} 的空间分辨率」；正文印「在实际催化体系中其显示的成像最佳
空间分辨率仍限于几十纳米，不足以实现化学分辨单个反应物或产物分子」——本分片的"观察"项据此
重建，与原书一致。强度剖面的具体定义（是否为半高宽）原书\textbf{未写明}，本分片亦不作断定。
\end{infer}

\begin{dataflow}
\textbf{从逐点采谱到一张可信的化学像，中间有四次信息转换}，每一次都可能改变结论。\\
（1）\emph{采谱}：针尖停在像素 $(x,y)$，积分 $t$ 秒，得到一条含背底的近场谱；\\
（2）\emph{减背底}：用同一针尖回撤或远离表面（源书用远离 5 nm 的谱线）作参考，扣除远场
与针尖自身的背景；\\
（3）\emph{取特征量}：把某条特征峰的峰高或峰面积作为该像素的"化学强度"。这一步丢掉了其余
全部谱形信息——若反应物与产物的谱带相互重叠，峰面积如何分解会直接改变化学像的形貌；\\
（4）\emph{拼像与配准}：把特征量按像素拼成化学像，再与同时记录的形貌像严格配准，才能说
"哪个位置是产物、哪个位置是反应物"。\\
因此，化学像的结论强度由四件事决定：像素间距是否小于空间分辨率（源书图 12-30（d）用剖面给出
20 nm，像素间距必须相应更小）、每像素积分时间是否足够、特征峰选取是否唯一、形貌与化学像是否
配准。缺任何一条，"某些位点活性更高"就只是候选解释而不是结论
\src{S2 §12.6.2，书内 p.932、p.934 / PDF p.418、p.420}。
\end{dataflow}

\paragraph{超高真空中的亚纳米极限}
在突破 TERS 成像分辨率极限方面，侯建国与董振超等率先实现了亚纳米尺度的 TERS，提出分辨
单个分子内部结构的光谱成像方法：把超高真空低温 STM（工作温度 78 K）与拉曼光谱结合，用 Ag
针尖在单晶 Ag(111) 基底上对一个卟啉类分子［H$_2$TBPP］内部的光谱特征成像，并确定了这些
特征如何随分子取向变化。其关键在于\textbf{通过调节 STM 探针尖端与基底之间的隧穿间隙来精确
调谐纳米腔的宽频等离激元共振频段}，使 LSP 共振谱与 Raman 跃迁相匹配，实现拉曼激发与拉曼
发射的双共振增强，从而获得超高分辨率的 TERS 成像（源书图 12-31、图 12-32）
\src{S2 §12.6.2，书内 p.932—933 / PDF p.418—419}。单分子成像的结果显示：分子外部的花瓣状
区域与分子中心区域给出不同的拉曼特征峰，且花瓣区的特征峰强度强于中心；用 1210 cm$^{-1}$ 及
以下波数的特征峰成像时，四个花瓣状结构清晰可见、与中心的对比度强，而用 1210 cm$^{-1}$ 以上
的特征峰成像时花瓣结构变模糊——说明低波数振动模式主要分布在花瓣区、高波数模式主要在分子
中心，与理论模拟相符。作者估计其 TERS 成像分辨率约为 0.5 nm，是迄今空间分辨率最高的分子
光谱成像\src{S2 §12.6.2，书内 p.933—934 / PDF p.419—420}。

这一结果打破了"受光可穿透金属针尖端部的深度所限，s-SNOM 与 TERS 在可见光和红外波段的空间
成像分辨率极限应在 10 nm 量级"的传统认识。源书给出的解释是：光场几乎可以限制在任意小的
区域，只受"所通过的均匀介质中电子不再表现为自由粒子的介质尺寸"限制；按 Thomas-Fermi 屏蔽
长度，这个尺寸界限约是 0.1 nm，低于该值则电子离域效应的影响变得显著
\src{S2 §12.6.2，书内 p.933 / PDF p.419}。同时要记住该实验的特殊性：与常规拉曼或 TERS 不同，
这里测得的 TERS 信号随入射激光功率增加呈\textbf{非线性增大}，作者将其归因于高阶非线性响应
产生的信号，且 TERS 信号对针尖的光学性质也表现出一定敏感度
\src{S2 §12.6.2，书内 p.934 / PDF p.420}。这意味着该条件下的信号强度不能简单按线性关系
换算成分子数或增强因子。

\paragraph{液相电化学体系中的 3 nm 成像}
除超高真空中的亚纳米结果外，任斌等展示了 TERS 成像在液相电化学体系中可获得 3 nm 的空间
分辨率：以单层 Pd 覆盖的单晶 Au(111) 表面为模型催化剂，通过电化学欠电位沉积在 Au(111) 上
构筑亚单层 Pd 与 Au 衬底的界面，以异腈苯（PIC）为探针分子，用 TERS 成像同时获得相关联的
STM 形貌与 TERS 光谱。结果是：与吸附在 Pd 平台位（terrace）的 PIC 相比，吸附在台阶位
（step edge）的 PIC 分子 N$\equiv$C 三键被削弱、振动频率降低，也导致其在台阶位更容易被
氧化（源书图 12-33）。该工作还发现台阶位的 TERS 信号强于平台位，理论计算表明可能源于台阶位
较强的等离激元光电场\src{S2 §12.6.2，书内 p.934—936 / PDF p.420—422}。其方法论价值在于：
TERS 在\textbf{实际（液相、电化学）条件下}也能实现高空间分辨的化学成像，并能在实空间解析
催化剂表界面的构效关系——但要注意，这里"台阶位更容易氧化"是靠\textbf{频率位移}（成键削弱）
与\textbf{强度变化}两条独立证据推出的，不是只看强度。

\begin{pitfall}
近场光谱的四类常见误判，共性是"把不属于近场的信息当成了近场信息"。\\
\textbf{（1）把远场信号当近场信号。}最典型的形式是针尖并未真正靠近就采谱：若针尖悬在几十
纳米之外、或根本没有实现焦斑—针尖—样品三者对准，读到的谱其实来自远场照明区（微米量级），
却会被当成纳米分辨的结果。\emph{判据}：必须做\textbf{针尖回撤对照}——源书图 12-31（c）中
针尖远离表面 5 nm 的黑线就是这种对照，近场特征峰在回撤后应大幅下降；若谱形几乎不变，或
化学像的分辨率不随针尖更换而改变，则读到的不是近场信号。\\
\textbf{（2）针尖污染造成的假谱带。}针尖末端吸附的无定形碳、有机物或样品残留会给出与样品
无关的谱带，且因为处于增强最强的位置，其强度往往很高。\emph{判据}：换新针尖后该谱带消失或
改变；该谱带在针尖远离表面时依然存在；或在同一样品的不同区域"到处都有、位置不动"。针尖
污染还会同时改变增强与分辨率，因此任何"新谱带出现"的结论都应先排除针尖因素。\\
\textbf{（3）把形貌像上的凸起直接当化学信息。}AFM／STM 形貌像上的凸起只说明几何或电子
结构有变化，与"是不是目标产物"无关；把（a）面板的高度起伏直接读成反应分布，等于跳过了
化学像。源书图 12-30 的做法是正确范式：先用 AFM 高度图（a）定位粒子，再用 1142 cm$^{-1}$
的特征峰做出化学像（b），最后用位点谱（c）与强度剖面（d）定量。反过来，若化学像的衬度与
形貌像一一对应、没有任何"化学上意外"的地方，也应怀疑是形貌串扰（针尖—样品距离变化引起
的信号变化）而不是化学衬度。\\
\textbf{（4）探针磨损后分辨率与增强同时下降而未被察觉。}针尖曲率半径变大或末端钝化，会使
增强因子下降、空间分辨率变差；表现在数据上是"信号变弱、像变模糊、峰位看似漂移"，很容易被
误读成样品变化、催化剂失活或反应衰减。\emph{判据}：在扫描前后用同一标准样品（如已知间距的
光栅或金纳米粒子）校验分辨率与信号强度；同一区域重复成像；发现"随时间单调变差"时首先怀疑
针尖而不是样品。\\
\textbf{（5）强激光引发的副反应被当成催化反应。}源书明确指出传统 TERS 常需强入射激光，可能
导致扩散、脱附甚至分子损坏\src{S2 §12.6.2，书内 p.931—932 / PDF p.417—418}。\emph{判据}：
降低功率后产物峰应仍能出现（只是信噪比下降）；或如源书那样用不引发反应的波长（633 nm）
作监测、用另一波长（532 nm）触发反应，把"光造成的"与"催化造成的"分开
\src{S2 §12.6.2，书内 p.929—930 / PDF p.415—416}。
\end{pitfall}

\subsubsection{纳米傅里叶变换红外光谱技术（Nano-FTIR）}
\label{ssec:nanoftir}

红外光谱是多相催化中广泛应用的表面分析技术，适应原位反应条件下的实时观测。红外辐射
（常用 2.5—25 µm 波长，对应 400—4000 cm$^{-1}$）可以激发分子振动、声子、表面等离激元以及
非金属导体的电子；通过对内禀振动与电子模式的直接测量，傅里叶变换红外光谱（FTIR）可提供
分子与固体材料的无标记"指纹"化学表征。然而红外光的微米级波长从根本上限制了 FTIR 用于多相
催化的纳米空间分辨成像；此外，即使在微米级分辨率下，要取得高信噪比的傅里叶变换图谱也必须
依靠高辐照度光源。亚波长成像在一定程度上可通过点扩散函数去卷积与衰减全反射技术实现，而
s-SNOM 的发展才使 FTIR 开始向纳米级空间分辨成像迈进
\src{S2 §12.6.3，书内 p.936 / PDF p.422}。

\begin{principle}
\textbf{原理：用针尖散射近场，用高次解调压制背底。}结合 FTIR 与 s-SNOM 的技术通常称为
Nano-FTIR：把红外光束聚焦到扫描探针（通常是 AFM）的金属针尖末端，利用所激发的 LSP 局部
增强散射的红外信号，在探针扫描样品表面的同时记录散射红外光的信息，进而产生表面的近场红外
图像。与 TERS 只需单色激光、可用 CCD 检测相比，Nano-FTIR 要在局域光谱信息的获得与光源
两方面提出更高要求。关键难点在于：\textbf{从 AFM 探针悬臂与针尖散射出来的远场背景信号远强于
从针尖末端散射出来的近场信号}，因此必须对近场红外信号的光强进行调制，才能从强散射背景中
提取近场信息。办法是利用针尖与样品之间纳米腔的强电磁场——通过针尖在 $z$ 方向的振动来调制
电磁场强，进而周期性调制散射的近场信号强度：当针尖以固定机械振动频率 $\Omega$ 振动时，
散射的近场红外信号在探测器中产生频率为 $n\Omega$（$n=1,2,3,4,\dots$）的周期性调制信号，
通过锁相放大处理即可提取这些近场信号，同时剔除未受调制的背景信号（直流部分），即
式~\eqref{eq:demod}。在一级近似上，Nano-FTIR 所得近场红外光谱信号的空间分辨率与探针针尖
顶点的\textbf{曲率半径}相当，从而不再受红外光波长限制
\src{S2 §12.6.3，书内 p.936—937 / PDF p.422—423}。由于需要调制，目前的 Nano-FTIR 常把 AFM
在轻敲模式（tapping mode）下与 FTIR 光谱仪联用。
\end{principle}

\begin{measure}
\textbf{光源是 Nano-FTIR 目前最重要的限制}，三类光源的取舍必须写清。\\
（1）\emph{传统热源}（加热灯丝或炽棒）：总光通量大，但缺乏空间相干性、难以聚焦，导致光谱
辐照度低，严重限制局域信号强度。历史上也尝试过用加热的针尖或样品直接做热近场成像，利用增强
的局域态密度与针尖散射光的空间相干性提高强度，但仍不足以实现高空间分辨红外成像，且容易
引起样品熔化等副作用。传统热源用于 s-SNOM 时只有强声子与等离激元可以被成像，较弱的分子
振动难以分辨\src{S2 §12.6.3，书内 p.937 / PDF p.423}。\\
（2）\emph{激光器}（商业化系统的常用选择）：提供高空间分辨率所必需的强频谱辐射。窄带连续波
（CW）可调谐激光器提供最高的光谱辐照度，但窄调谐范围使所采纳米光谱通常只能覆盖 1 个振动
模式；飞秒脉冲红外激光器增加了光谱带宽、可覆盖多个振动模式，却牺牲了光谱辐照度与功率，
难以实现覆盖整个中红外区且具备足够功率的高重频光源。总体而言，目前最好的红外激光源带宽也只
是常规红外热源范围的一小部分；对连续波与飞秒脉冲光源而言，通过光谱调谐覆盖更宽范围通常
伴生激光与样品的漂移，限制光谱质量\src{S2 §12.6.3，书内 p.937 / PDF p.423}。\\
（3）\emph{同步辐射红外光源}：宽频谱、高亮度、空间相干，克服了上述限制——光谱辐照度是热源的
1000 倍，光谱带宽比最宽的激光源还宽 10 倍以上；其空间相干性与高光谱辐照度使光可以按衍射
限制聚焦，从而提高入射光与 AFM 针尖的耦合；光源稳定性（采集期间光照波动小、功率与光谱分布
稳定）带来更好的信噪比，使有效信号可在几分钟到几十分钟内取得。源书的判断是：在光谱辐照度、
光谱范围与重复频率的综合权衡下，同步辐射红外是目前可在单次采集中获得整个中红外范围近场
光谱的最优来源。结合同步辐射红外光源的 s-SNOM 已展示以优于 40 nm 的空间分辨率获得覆盖整个
中红外指纹区（700—5000 cm$^{-1}$）的分子振动光谱图，采集时间约几秒，基于此的化学敏感红外
光谱成像可在几分钟内收集完；这种方法通常称为 SINS（synchrotron-radiation-based infrared
nanospectroscopy），其检测灵敏度为数百个分子或化学键
\src{S2 §12.6.3，书内 p.937—938 / PDF p.423—424}。
\end{measure}

\paragraph{催化实例：Pt 颗粒上 NHC 分子的可逆氧化还原}
Gross 与 Toste 等用 SINS 研究了直径约 100 nm 的 Pt 颗粒表面分子的反应性：把 N-杂环卡宾
分子（NHCs）锚定在 Pt 纳米颗粒表面、沉积在硅片上，再用 SINS 研究这些结合在 Pt 颗粒上的
NHC 分子的化学转化，光谱成像空间分辨率可达 25 nm（源书图 12-34）。结果是：氧化条件下醇
羟基（—OH）可被氧化为羧酸基团（—COOH），还原条件下羧酸基团可逆地还原为醇羟基。利用 SINS
记录粒子不同位点处的官能团信息，即可得到相应位点的催化活性，进而分辨具有不同活性的颗粒
区域：\textbf{与颗粒顶部的平坦区域相比，包含低配位数金属原子的颗粒边缘催化活性更高}，
能更有效地催化 NHC 分子中化学活性基团的氧化与还原（源书图 12-35）。此外，样品暴露于轻度
氧化/还原条件与苛刻条件下的 IR 光谱相似，说明锚定分子在 Pt 颗粒上的不同取向对吸附性质影响
不大，除非发生化学变化导致本征基团性质改变\src{S2 §12.6.3，书内 p.938—939 /
PDF p.424—425}。

这项研究把红外光谱的空间分辨率比远场技术提高了三个数量级，但也暴露了当前的折中：所研究的
颗粒（约 100 nm）仍比实际催化剂 1—10 nm 的尺寸范围大，原因在于红外纳米光谱的空间分辨率
受针尖尖端直径（约 20 nm）限制。Nano-FTIR 当然也能检测小纳米颗粒（$<10$ nm）上表面分子的
信号，但此时只能给出其表面所有吸附分子的\textbf{平均}红外信号，无法区分不同吸附部位的分子
\src{S2 §12.6.3，书内 p.939 / PDF p.425}。

\paragraph{时间分辨率的现实}
即便利用同步辐射红外光源，目前每个单点的 Nano-FTIR 谱图测量仍需几分钟，因此无法原位实时
监测表面催化化学反应——反应的典型持续时间（$<1$ s）比单点红外光谱的获取时间短得多。在
Gross 等的研究中，取得单点 FTIR 谱图约需 20 min，导致他们无法实现 Nano-FTIR 成像、仅限于
线扫描测量；要对更大面积成像在技术上更为挑战，还会引起长时间扫描的热漂移等问题
\src{S2 §12.6.3，书内 p.939—940 / PDF p.425—426}。注意这两个时间尺度（SINS 单谱约几秒、
成像几分钟；Gross 等单点约 20 min）并不矛盾：它们对应不同的装置条件与信噪比要求，说明
\textbf{近场红外的时间分辨率强烈依赖具体系统}，报告时必须给出采谱时间与像素数，否则"实时"
二字无法比较。

\paragraph{偶极选律：哪些振动模式根本看不见}
在 Nano-FTIR 中，针尖极化了沿针尖长轴方向的电场，因此测量\textbf{主要对偶极矩平行于针尖的
振动模式敏感；偶极矩平行于表面的振动将不被检测到}——这与远场单晶反射红外的表面选律相同。
在已展示的研究中，对表面单层分子的振动，目前比较容易检测具备高红外吸收截面的键，如
C=O 与 O—H 伸缩模式；而 C—H 与 C—C 伸缩模式的检测尚较具挑战性，需要利用可激发表面等离
激元场的金属（如 Au）来增强局域 IR 信号\src{S2 §12.6.3，书内 p.940 / PDF p.426}。这条选律
是判读 Nano-FTIR 谱图时最容易出错的地方：\textbf{某个峰"没有出现"不等于该基团不存在}，
它可能只是取向不对。

\begin{crosscheck}
\textbf{Nano-FTIR 与 TERS 是互补技术，不是替代关系}——两者都通过针尖与光的相互作用克服衍射
极限、以高空间分辨率检测表面分子振动，但权衡相反。Nano-FTIR 的优势有两条：① 红外吸收截面
比拉曼高很多，所需光通量较低，减少了光诱导副反应的可能；② 与拉曼相比，IR 散射或吸收是
\textbf{线性现象}，因此限域增强红外信号受针尖结构变化的影响较小。其代价是需要高亮度、高
相干的宽带光源（同步辐射），且较长采谱时间导致时间分辨率有限；TERS 则可以在样品表面以几秒
的时间分辨率识别吸附物种、跟踪化学反应\src{S2 §12.6.3，书内 p.940 / PDF p.426}。推论很
直接：\textbf{要跟踪反应的时间演化用 TERS，要做广谱指纹识别且怕光损伤用 Nano-FTIR}；若
两者结论不一致，先比较它们的选律（TERS 无偶极取向限制的同类问题，但受共振条件与针尖增强
体积影响）与采谱时间，而不是先怀疑样品。\\
此外，源书对 Nano-FTIR 的分辨率上限给出了明确的保留：TERS 在超高真空中展现的亚纳米分辨率
说明现有 Nano-FTIR 在空间分辨率上仍有很大提高空间，而\textbf{不能简单地把扫描近场光学显微镜
的空间分辨率归结为由针尖顶端直径决定}。目前商业化 s-SNOM 系统为减少针尖对样品的破坏与
干扰，通常采用轻敲模式 AFM 以最小化针尖—表面相互作用、便于检测弱结合物种，但这同时造成
针尖与表面之间距离较长，减少了红外信号的增强；s-SNOM 装置的稳定性、振动隔离与新型针尖的
开发都有望进一步改善其灵敏度与空间分辨率
\src{S2 §12.6.3，书内 p.940 / PDF p.426}。
\end{crosscheck}

\subsubsection{光学激发的针尖探测方法}
\label{ssec:afmir}

另一类利用扫描探针实现纳米级光谱与亚波长成像的思路是：\textbf{用光学手段激发，而用扫描探针
探测与成像}。这类方法可与不同波段的光结合，以克服 s-SNOM 的一些弱点，其中发展较快的是基于
AFM 的红外光谱技术（atomic force microscopy-based infrared spectroscopy，AFM-IR）
\src{S2 §12.6.4，书内 p.940 / PDF p.426}。

\begin{principle}
\textbf{AFM-IR 与 s-SNOM 的根本差别：测吸收还是测散射。}AFM-IR 测量样品\textbf{吸收}的光，
而 s-SNOM 测量样品\textbf{散射}的光。因此：样品是强烈光散射体时，s-SNOM 表现较佳；而对较弱
的红外散射、且针尖与样品作用较复杂的情况，AFM-IR 可以相对快速、直接地获得吸收光谱
\src{S2 §12.6.4，书内 p.940—941 / PDF p.426—427}。\\
\textbf{工作原理}：用 AFM 悬臂针尖检测由 IR 辐射引起的局部吸收所产生的样品热膨胀，从而获得
局域红外吸收光谱。实验室设备中，AFM-IR 通常用可调谐红外激光器聚焦照射探针尖端附近的样品
区域，IR 吸收引起对悬臂尖端探测原子力的冲击，造成悬臂振荡；测量悬臂对 IR 吸收的响应随波长
的变化，即可在纳米尺度得到该表面区域的 IR 吸收光谱。目前对 AFM 针尖所探测的瞬态力机制尚不
十分明晰，较通用的理解是：红外辐照使样品局域温度升高，带来瞬时热膨胀，热膨胀对悬臂振荡
给出一个冲击力。大多数情况下，\textbf{悬臂的振荡幅度与吸收的光量成正比，也与样品的吸收系数
成正比}；这些性质在样品的特定点保持恒定，因此在固定点测量悬臂振荡幅度随波数的变化，即可
提供与吸收系数成正比的信号，即局域 IR 吸收光谱
\src{S2 §12.6.4，书内 p.941 / PDF p.427}：
\begin{equation}\label{eq:afmir}
A_{\mathrm{cant}}(\lambda)\;\propto\;\mu_{\mathrm{abs}}(\lambda)\,I(\lambda)\,\chi_{\mathrm{th}}
\end{equation}
式中 $A_{\mathrm{cant}}$ 为悬臂峰—峰（或本征模）振荡幅度，$\mu_{\mathrm{abs}}$ 为样品吸收
系数，$I$ 为局域光强，$\chi_{\mathrm{th}}$ 为把吸收能量转换为悬臂响应的热—力转换系数。
\emph{核对说明}：已对照原页面图像（书内 p.941 / PDF p.427）。原书只作文字说明：「在大多数
情况下，悬臂的振荡幅度与吸收光的量成正比，也与样品的吸收系数成正比」，并明确「目前，我们对
AFM 针尖所探测的由于红外吸收导致的瞬态力的机制尚不太明晰」——原书\textbf{未给出}表达式，也
\textbf{未定义}热—力转换系数。式~\eqref{eq:afmir} 按该段文字补写，其中 $\chi_{\mathrm{th}}$
是本分片为把比例关系写完整而引入的记号，故本条只能作定性关系使用。
\end{principle}

\begin{measure}
AFM-IR 的测量与边界条件。数据采集通常通过测量悬臂振荡幅度随波长（光谱）或样品位置
（化学成像）的变化进行，主要关注峰—峰振荡幅度或悬臂本征模振荡幅度；在某些情况下可采用
高阶模式提高信噪比，并改善由靠近悬臂的空气被加热所产生的\textbf{非局部背景力}的排斥。采用
高重复频率的 QCL 激光器匹配 AFM 悬臂共振，可以把悬臂振荡从衰减的瞬态振荡变为连续波振荡，
从而允许检测更多的红外吸收、显著提高检测灵敏度。但相对 s-SNOM 这一表面敏感技术，AFM-IR 的
红外吸收深入到样品更深处，目前主要适用于\textbf{高厚度的分子聚集体}（$>40$—80 nm），
尚未出现可分辨单层或亚单层吸附物种结构的催化或表面反应报道；这一对表界面单分子层不敏感的
特点，限制了它目前在催化反应中的应用\src{S2 §12.6.4，书内 p.941 / PDF p.427}。
\end{measure}

\paragraph{IRSTM：把红外吸收搬到隧道结里}
光学激发的针尖探测方法的空间分辨率极限目前仍有很大提高空间。正如 TERS 所展示的，超高真空
STM 的亚原子级空间分辨与更精细的调制可以提高光谱成像的空间分辨率与表面灵敏度；同样地，
基于超高真空 STM 开发的红外扫描隧道显微镜（IRSTM）已可表征亚单层分子在固体表面的红外吸收，
表现出比 IETS 更好的分子振动谱分辨率，并有研究用它区分精细的化学结构，如单晶 Au(111) 表面
的两种四金刚烷异构体（[121]tetramantane 与 [123]tetramantane）
\src{S2 §12.6.4，书内 p.941 / PDF p.427}。其主要原理是：在可变波数的单色红外辐照下，用
隧穿模式下的 STM 针尖检测分子红外吸收导致的样品热膨胀。表面没有分子吸收时，恒电流模式下的
STM 针尖检测到晶体样品本身沿垂直表面方向的热膨胀；表面存在吸附分子时，每当波长被调谐到
分子振动共振处即出现分子吸收、激发分子振动态，该激发态通常在 10 ps 内弛豫并把能量转移到
基底上，从而在该波长处引起\textbf{额外的表面膨胀}。判据就是"有无分子之间的表面膨胀差异"。
目前的研究表明，IRSTM 获得的红外光谱与传统的远场红外光谱可以相互对应
\src{S2 §12.6.4，书内 p.941—942 / PDF p.427—428}——这为近场谱的峰位归属提供了一个独立的
回查渠道。

\paragraph{SX-STM：X 射线 + 隧道结，以及它的三重困难}
除分子振动光谱外，光学激发的针尖探测方法还可帮助其他表面化学分析手段提高空间分辨率，其中
一个充满希望的结合是超高真空 STM 与同步辐射 X 射线，即同步加速器 X 射线扫描隧道显微镜
（SX-STM）：在隧穿模式下用单色 X 射线照射针尖的隧穿结，以 X 射线激发调制隧穿电流；当 X 射线
能量处于样品局域元素的共振吸收位置时，隧穿电流会大幅增大，从而获得化学衬度信息
\src{S2 §12.6.4，书内 p.942 / PDF p.428}。但实现亚纳米 X 射线吸收谱还有三重困难：第一，
X 射线激发的电子不仅穿入针尖尖端，也可由针尖\textbf{侧壁}吸收，这降低了空间分辨率；第二，
从样品中喷出的光电子在远场提供了一种\textbf{欺骗性的电流信号}，STM 反馈控制器会把针尖提出
隧穿状态；第三，测得的针尖电流是常规隧穿电流（与表面起伏相关）与 X 射线激发电子（携带
材料化学信息）的\textbf{总和}，必须设法提取后者。通过开发新型智能针尖与表面形貌过滤器，
目前已在解决上述问题上取得进展，在一定程度上展现了具有原子灵敏度的直接元素成像的可能；
但这些技术仍处在概念性开发阶段，要用于实际催化研究还有很长的路要走
\src{S2 §12.6.4，书内 p.942 / PDF p.428}。

\subsection{结束语}
\label{sec:nano-end}

\paragraph{三种技术的分工}
把本节的三种纳米光谱技术放在同一张表里对照，差异一目了然——它们测的物理量不同，因而分辨率
来源、样品要求和主要假象都不同。

\begin{table}[htbp]
\centering
{\srcsize
\setlength{\tabcolsep}{4pt}
\begin{tabularx}{\textwidth}{@{}p{15mm}XXXX@{}}
\toprule
技术 & 测什么（物理量） & 分辨率由什么决定（源书报道值） & 样品与环境要求 & 主要假象 \\
\midrule
\textbf{TERS} &
针尖正下方极小体积内分子的\textbf{拉曼散射}（振动指纹），与形貌同时得到 &
针尖末端 LSP 增强场的体积，与曲率半径同量级；大气/液相 20 nm、液相电化学 3 nm、超高真空低温 0.5 nm &
金属针尖（Ag 或 Ag/Au 包覆）；分子须可吸附且有拉曼活性；液相或超高真空/低温；单层可测 &
针尖污染产生假谱带；针尖未真正靠近时读到远场信号；强激光致扩散、脱附、副反应；针尖磨损后
增强与分辨率同时下降 \\
\textbf{Nano-FTIR}（s-SNOM + FTIR） &
针尖\textbf{散射}的红外近场信号，即局域红外吸收谱 &
针尖顶点曲率半径（约 20 nm）；常规 s-SNOM 40 nm 量级；SINS 在 Pt 颗粒上 25 nm &
需要高亮度、空间相干的宽带红外光源（激光器或同步辐射）；AFM 轻敲模式；弱散射样品更有利 &
远场背底未压住（悬臂与针尖杆散射）；偶极选律使部分振动模式"消失"（C—H、C—C 难检）；
长时间扫描的热漂移 \\
\textbf{AFM-IR} &
样品\textbf{吸收}红外后的热膨胀对悬臂的冲击，悬臂振幅正比于吸收系数 &
针尖下方受热膨胀的体积（受热扩散加宽）；源书未给出报道值，并指出该技术对单层不敏感 &
可调谐红外激光器（QCL 与悬臂共振匹配）；样品需较厚（$>40$—80 nm 的分子聚集体） &
空气受热产生的非局部背景力；厚度与热学性质衬度被当作化学衬度；瞬态力机制本身尚不明晰 \\
\bottomrule
\end{tabularx}
}
\caption{三种纳米光谱技术的分工。读者应横着读：同一行里"测什么"决定了"分辨率由什么决定"，
也决定了"样品必须满足什么"与"最可能出现哪种假象"。三者的分辨率来源都不是波长：TERS 与
Nano-FTIR 靠针尖末端的局域增强体积，AFM-IR 靠针尖下方受热膨胀的体积。报道值均取自源书正文，
书内 p.929—942 / PDF p.415—428；AFM-IR 一列源书未给数值，仅按原理描述。
\emph{核对说明}：已对照原页面图像（书内 p.929—942 / PDF p.415—428）。本表是跨小节的归纳，
原书\textbf{并无此表}；各格内容均可在正文找到出处（TERS 与 Nano-FTIR 的分辨率由针尖末端的局域
增强体积决定、AFM-IR 由针尖下方的热膨胀体积决定，分别见书内 p.929、p.936、p.941）。"主要假象"
一列中「厚度与热学性质衬度被当作化学衬度」一条系依原理所作的推论，原书\textbf{未逐条给出}
实验判据，列出以便使用时自查。\src{据 S2 §12.6.2—§12.6.4，书内 p.929—942 / PDF p.415—428 归纳。}}
\label{tab:nano-three}
\end{table}

\paragraph{方法一览与各自的能力边界}
下表把本节涉及的全部方法按"激发/探测方式—报道分辨率—关键限制"排成一张长表，便于查用。

{\srcsize
\setlength{\tabcolsep}{3pt}
\begin{longtable}{@{}p{24mm}p{42mm}p{32mm}p{48mm}@{}}
\caption{近场与纳米光谱方法一览。读者应重点看最后一列：每一项技术的"能不能用"都由该列的
限制决定，而不是由分辨率数字决定。表中数值均取自源书正文（书内 p.926—942 / PDF p.412—
428）；凡源书未给出数值者如实标注。\src{据 S2 §12.6—§12.6.4，书内 p.926—942 / PDF
p.412—428 归纳。}}
\label{tab:nano-methods}\\
\toprule
方法 & 激发／探测方式 & 空间分辨率（源书报道值） & 关键限制与应用边界 \\
\midrule
\endfirsthead
\toprule
方法 & 激发／探测方式 & 空间分辨率（源书报道值） & 关键限制与应用边界 \\
\midrule
\endhead
\bottomrule
\endlastfoot
孔径式 NSOM & 金属（Al／Ag）涂覆锥形光纤，孔径约 100 nm；透射、反射、荧光或收集模式 &
由孔径尺寸决定（亚波长；源书未给具体数值） &
可见光透过率仅 $10^{-7}$—$10^{-4}$（几十纳瓦）；中红外 50 nm 孔径降至 $10^{-30}$ 量级，
红外近场成像几乎不可行 \\
PSTM／STOM & 样品背面全内反射产生隐失波，光纤探针端部检测隐失场并转为传播模式 &
由探针端部与隐失场衰减长度决定（源书未给数值） &
近场光学像受样品调制，但图像解析较复杂；源书称目前尚未见用于催化研究的报道 \\
散射式 s-SNOM & 激光聚焦照射金属针尖，收集散射光，以调制—高次解调提取近场信号 &
仅由针尖曲率半径决定；可能优于 10 nm &
远场背底远强于近场信号，必须有效抑制；分辨率的传统认识上限为 10 nm 量级（已被 TERS 打破） \\
TERS（超高真空低温 STM） & 激光 + 隧穿控制的 STM（78 K，Ag 针尖，Ag(111)，H$_2$TBPP 分子） &
0.5 nm（单个分子内部的光谱成像） &
需超高真空与低温；信号随激光功率呈非线性增大；对针尖的光学性质敏感；靠隧穿间隙调谐 LSP 共振 \\
TERS（液相电化学） & 电化学 STM + 拉曼，Pd/Au(111) 模型催化剂，PIC 探针分子 &
3 nm &
台阶位与平台位的信号强度本身不同（等离激元光电场差异）；归属依赖频率位移 + 强度两条证据 \\
TERS（大气／模型体系） & 双波长 532／633 nm；Ag 包覆或 Al$_2$O$_3$ 保护针尖 &
20 nm（Ag 纳米粒子上的 DMAB 光催化成像） &
强激光可致扩散、脱附、分子损坏；单波长激发会混淆"自然过程"与"光催化反应" \\
Nano-FTIR（实验室光源） & 可调谐 CW 或飞秒红外激光 + AFM 轻敲模式 + FTIR 干涉 &
与针尖曲率半径相当（源书示例约 20 nm） &
窄带光源通常只覆盖 1 个振动模式；宽带飞秒源牺牲辐照度与功率；调谐时伴生漂移 \\
SINS（同步辐射红外） & 同步辐射红外 + s-SNOM；光谱辐照度为热源的 1000 倍，带宽比最宽激光源宽 10 倍以上 &
优于 40 nm；Pt 颗粒上 25 nm &
依赖同步辐射装置；单点采谱与成像时间随信噪比要求可在几秒到 20 min 之间变化 \\
AFM-IR & 可调谐红外／QCL 激发，悬臂检测样品热膨胀（测吸收，非测散射） &
针尖下方受热膨胀体积（源书未给数值） &
仅适用于较厚分子聚集体（$>40$—80 nm）；对单层／亚单层吸附物种不敏感；非局部背景力 \\
IRSTM & 超高真空 STM + 可变波数单色红外，检测分子吸收引起的表面热膨胀 &
亚单层分子红外吸收；可区分 Au(111) 上的四金刚烷异构体 &
需超高真空；依赖分子振动态约 10 ps 弛豫引起的额外表面膨胀；谱与远场红外可相互对应 \\
SX-STM & 单色 X 射线照射隧穿结，检测共振吸收位置处隧穿电流的增大 &
目标为原子级元素成像（概念性开发阶段） &
三重困难：针尖侧壁吸收 X 射线激发电子、光电子远场假电流使反馈提出隧穿态、隧穿电流与激发
电子电流混合需分离 \\
\end{longtable}
}

\paragraph{共同陷阱}
上面表格最后一列汇总了各技术特有的假象；另有四条误判是\textbf{跨技术共有}的，判断顺序建议
固定下来。\\
（1）\emph{没有回撤对照就把谱当近场谱}：任何近场实验至少要有一条"针尖离开表面"的对照谱。
（2）\emph{不查针尖状态就把新谱带当新物种}：换针尖、重复扫描、在同一区域多处采谱，三条中
至少做一条；针尖污染的特征是"到处都有、位置不动、换针即消"。\\
（3）\emph{把形貌衬度当化学衬度}：化学像必须独立于形貌像成立——若两者一一对应且毫无意外，
先怀疑是针尖—样品距离串扰。\\
（4）\emph{不监控针尖磨损就解释信号衰减}：增强与分辨率同时下降是针尖钝化的签名，不是
催化剂失活的签名；长扫描前后用标准样品校验一次，成本远低于一次错误结论。

\begin{crosscheck}
近场结论的独立回查途径，按成本从低到高排列：\textbf{① 峰位对照}——把近场谱峰位与常规
远场拉曼／红外谱对照（源书明确说明 IRSTM 的红外光谱与远场红外光谱可以相互对应），峰位一致
而相对强度不同，通常是取向或选律效应，而不是新物种；\textbf{② 距离依赖}——近场信号应随
针尖—样品间距按式~\eqref{eq:nearfield-decay} 迅速下降，远场信号不会，这是区分二者最直接的
实验；\textbf{③ 换针尖／换位点重复}——用于排除针尖污染与个例；\textbf{④ 与形貌和粒径分布
配准}——把化学像与 AFM 高度像、粒径统计逐粒对照，检验"边缘／台阶位活性更高"这类构效关系
是否稳健（源书 SINS 与 TERS 两处结论都建立在这种配准之上）；\textbf{⑤ 与理论计算对照}——
源书在 H$_2$TBPP 单分子成像（花瓣区与中心的振动模式分布）与 Pd 台阶位增强（等离激元光电场）
两处都用理论模拟作为独立支撑；\textbf{⑥ 与元素分析对照}——对组分归属存疑时，用电子显微与
元素分析给出独立组分信息，因为近场振动谱只回答"是什么基团"，不回答"含什么元素"。
最后要记住本节反复出现的边界：纳米光谱技术尚处起步阶段，还未能广泛用于多相催化剂的原位
实时研究；现有研究已表明其可行性，但一张近场化学像的定量结论，必须连同针尖、间距、解调
阶数、采谱时间与环境条件一起报告，否则无法复核
\src{S2 §12.7，书内 p.942 / PDF p.428}。
\end{crosscheck}

扫描隧道显微镜与原子力显微镜在过去的 30 年间取得了长足的技术进步，在催化与表面科学领域
发挥着重要作用：它们在催化剂表面结构的原子解析、表面催化反应机理与动力学的可视化方面有着
独到优势，且能在不同反应气氛、压力和介质环境下工作。而把扫描探针技术与分子光谱技术结合
以实现纳米光谱技术（如针尖增强拉曼光谱与纳米傅里叶变换红外光谱），有望实现催化剂原子结构
与反应／产物分子振动结构的联合表征与关联。可以看到，扫描探针技术的发展为纳米光谱技术的
开发提供了很好的机会；尽管纳米光谱技术尚处于起步阶段、还未能广泛用于多相催化剂的原位实时
研究，但现有研究已经表明了将其用于催化研究的可行性，并且还有很大的发展空间使之成为新一代
催化研究的必备工具\src{S2 §12.7，书内 p.942 / PDF p.428}。

\srcnote{本节（12.6—12.7）内容据 S2（下册）第 12 章，书内 p.926—942 / PDF p.412—428。
源书图 12-28 至图 12-35 为原书配图，本分片未复制页面图像，正文中对它们的描述均转述自图注
与相关段落；本分片的两幅图分别为据源文正文改绘的示意图与用于说明数量级关系的模型曲线，
均已在图题中标注。}
